\begin{tabularx}{\textwidth}{L{32mm}X X}
\toprule
\textbf{Fogalom} & \textbf{Jelentés} & \textbf{Vizsgán kiemelendő gondolat} \\
\midrule
Hoszt & A hálózat végpontja, amely alkalmazásokat futtat és adatot küld vagy fogad. & Például kliensgép, szerver, mobil eszköz vagy beágyazott eszköz. \\
Csomópont & Általános hálózati elem, amely több átviteli vonalhoz kapcsolódhat. & Lehet végpont vagy kapcsolóelem; feladata lehet a továbbítás és útválasztás. \\
Átviteli vonal / csatorna & Két vagy több hálózati elem közötti kommunikációs közeg. & Lehet vezetékes, optikai vagy rádiós; a fizikai réteghez kapcsolódik. \\
Protokoll & Szabályok és konvenciók halmaza két társelem adatcseréjéhez. & Meghatározza a szintaxist, szemantikát és az időzítési/sorrendi szabályokat. \\
Interfész & Két szomszédos réteg kapcsolódási pontja és szolgáltatásleírása. & A felső réteg szolgáltatást kér, az alsó réteg elrejti a megvalósítás részleteit. \\
Szolgáltatás & Amit egy réteg a fölötte lévő rétegnek nyújt. & A szolgáltatás a „mit”, a protokoll inkább a „hogyan” kérdésére válaszol. \\
Társelem & Különböző gépeken azonos rétegben lévő funkcionális elemek. & Látszólag egymással kommunikálnak, valójában az adat lefelé és felfelé halad a rétegek között. \\
SAP & Service Access Point, vagyis szolgálatelérési pont. & Itt érhető el egy réteg szolgáltatása; címmel azonosítható. \\
PDU & Protocol Data Unit, egy réteg protokoll-adategysége. & Az alsóbb réteg fejléccel és esetleg utótaggal látja el a felsőbb rétegből kapott adatot. \\
\bottomrule
\end{tabularx}
